\section{Relation with Actual Causality} \label{sec:relation_with_ac}

 Sections~\ref{sec:shap_examples}--\ref{sec:DCE} have
positioned \ourapproach against feature-attribution baselines. We now turn to
\ourapproach's relation to the causal-attribution tradition it inherits the
witness mechanism from --- Halpern's actual causality. This section establishes
the \emph{formal} connection (Theorems~\ref{th:ac-exp}--\ref{th:exp-ac} and
Proposition~\ref{prop:filter}); Section~\ref{sub:pearl_actual_cause_prob} adds
the side-by-side comparison with Pearl's \emph{probability of actual
causation}; and Section~\ref{sec:ac_benchmark} measures the empirical gap in
runtime and verdict recovery.

\ourapproach{} is not meant to explicate Halpern's notion of actual causality; the goal
is to combine the intuitions underlying Halpern's notion with those motivating Pearl's
probabilities of necessity and sufficiency to provide a more general tool for causal
attribution. Nevertheless, to show that we preserve the spirit of Halpern's definition,
we now record the connection between actual causality and the joint necessity-sufficiency
machinery of Section~\ref{sec:causal_impact}. The goal is to show that there is a
natural choice of variable selection distribution $\Gamma$ and alternative-value
distribution $\Delta$ under which \ourapproach{} verdicts agree with actual-causality
verdicts. With one caveat: probability-of-causation maximisation will not identify
subset-minimal causes that are not cardinality-minimal, a misalignment we
characterise and remedy below.

\medskip\noindent\textit{Scope: a necessity-only correspondence.} Halpern's definition
involves factivity, context-sensitive necessity, and minimality, but no sufficiency
clause analogous to Pearl's $\mathrm{PS}$ or $\mathrm{PNS}$. The correspondence we
establish in this section is therefore \emph{necessity-only}: the sufficiency component
of the kernel $\Phi$ defined below acts as a factivity check (restoring factual values
must preserve $\varphi$ under fixed noise), not as a substantive sufficiency-attribution
component. The substantive sufficiency component of \ourapproach{}, the $Y^s$ factor
in Definition~\ref{def:jointnecsuf} and the $ci_S$ term in
Section~\ref{sub:synthetic}, is exercised in the synthetic overdetermination and
undercutting evaluation (Section~\ref{sub:synthetic}), the differential causal effect
comparison (Section~\ref{sec:DCE}), and the SIR benchmark
(Section~\ref{sec:sir_benchmark}), all of which are settings where AC cannot adjudicate.
The narrowness of this section's correspondence reflects what AC itself permits, not a
limitation of \ourapproach{}.

We first recall Halpern's definition.

\begin{definition}[Actual Cause, after \citealp{actualCausalityHalpern}] \label{def:ac}
    \(\mathbf{C} = \mathbf{c}\) is an actual cause of \(\varphi\) in the causal setting \((M, \mathbf{u})\) if the following three conditions hold:
    \begin{itemize}
        \item \textbf{Factivity}: \((M, \mathbf{u}) \vDash (\mathbf{C} = \mathbf{c}\,)\) and \((M, \mathbf{u}) \vDash \varphi\),
        \item \textbf{Context-sensitive Necessity}: There is a set of variables \(\mathbf{T} \subseteq \mathbf{V}\) and an alternative setting \(\mathbf{c'}\) of variables in \(\mathbf{C}\) such that if \((M, \mathbf{u}) \vDash \mathbf{T} = \mathbf{t^\star}\), then
            \[(M, \mathbf{u}) \vDash [\mathbf{C} \leftarrow \mathbf{c'}, \mathbf{T} \leftarrow \mathbf{t^\star}] \neg \varphi,\]
        \item \textbf{Minimality}: \(\mathbf{C}\) is minimal; there is no strict subset \(\mathbf{C_{sub}} \subset \mathbf{C}\) such that \(\mathbf{C_{sub}} = \mathbf{c_{sub}}\) satisfies the above two conditions where $\mathbf{c_{sub}}$ is the restriction of \(\mathbf{c}\) to the variables in \(\mathbf{C_{sub}}\).
        \end{itemize}
\end{definition}

\noindent To connect this to \ourapproach{}, we instantiate the joint
necessity-sufficiency machinery (Definition~\ref{def:jointnecsuf}) at a single
configuration $(\mathbf{C}, \mathbf{T})$ and a fixed noise realisation $\mathbf{u}$,
and read off the configuration-level contribution. The resulting kernel combines the
same three ingredients as Halpern's definition: a sufficiency check at factual values,
a necessity check against alternatives, and a $\Gamma$-weight that records how much
the configuration is being asked about.

\begin{definition}[AC-aligned PCI kernel] \label{def:ackernel}
Fix a PSCM $M=\langle\mathbf{U},\mathbf{V},\mathbf{F},P_{\mathbf{U}}\rangle$,\footnote{%
The kernel and the theorems below take the full PSCM (structural equations $\mathbf{F}$
plus noise distribution $P_{\mathbf{U}}$) as given, since both the sufficiency and the
necessity factors require sampling counterfactual outcomes under $\mathbf{u}$. As
discussed in Section~\ref{sec:intro}, this is a stronger commitment than approaches
that bound $\mathrm{PN}/\mathrm{PS}/\mathrm{PNS}$ from observational data and a causal
graph alone, but is the price of pointwise (not just bound) verdicts at the
configuration level.} a noise
realisation $\mathbf{u}$, suspect and witness sets $\mathbf{S},\mathbf{W}\subseteq\mathbf{V}$
with factual values $\mathbf{s}^\star,\mathbf{w}^\star$, a variable selection
distribution $\Gamma$ on $2^{\mathbf{S}}\times 2^{\mathbf{W}}$ with mass function
$p^\Gamma$, and an alternative-value distribution $\Delta(\mathbf{s}^\star)$. For an
outcome event $\varphi$, active suspect set $\mathbf{C}\subseteq\mathbf{S}$ and active
witness set $\mathbf{T}\subseteq\mathbf{W}$ with $\mathbf{T}\cap\mathbf{C}=\emptyset$,
define
\begin{align}
\Phi(\mathbf{C},\mathbf{T})
&\;:=\;
p^\Gamma(\mathbf{C},\mathbf{T})
\;\cdot\;
\mathbb{I}\!\left[(M,\mathbf{u})\vDash
[\mathbf{C}\!\leftarrow\!\mathbf{s}^\star\!\mid_{\mathbf{C}},\,
 \mathbf{T}\!\leftarrow\!\mathbf{w}^\star\!\mid_{\mathbf{T}}]\,\varphi\right]
\nonumber \\
&\quad\;\cdot\;
\int
\mathbb{I}\!\left[(M,\mathbf{u})\vDash
[\mathbf{C}\!\leftarrow\!\mathbf{c}',\,
 \mathbf{T}\!\leftarrow\!\mathbf{w}^\star\!\mid_{\mathbf{T}}]\,\neg\varphi\right]
\Delta_{\mathbf{C}}(\mathbf{s}^\star)(d\mathbf{c}').
\label{eq:ackernel}
\end{align}
\end{definition}

\medskip\noindent Concretely, $\Phi(\mathbf{C},\mathbf{T})$ is the
matched-configuration slice of the joint necessity-sufficiency measure $P_k^{s,n}$
from Definition~\ref{def:jointnecsuf}: integration over $P_{\mathbf{U}}$ is replaced
by point evaluation at $\mathbf{u}$, and both the sufficiency and the necessity sums
are pinned to the same $(\mathbf{C},\mathbf{T})$ pair, mirroring AC's structure of
testing factivity and the alternative against a single suspect/witness pair. The
kernel is positive exactly when three conditions hold simultaneously:
$(\mathbf{C},\mathbf{T})$ has positive $\Gamma$-mass; restoring $\mathbf{C}$ to its
factual values while holding $\mathbf{T}$ fixed preserves $\varphi$ under $\mathbf{u}$
(the factivity-as-sufficiency component); and there exists a
$\Delta$-positive-density alternative $\mathbf{c}'$ for which the corresponding
intervention overturns $\varphi$ under $\mathbf{u}$ (the necessity component). The
first factor takes the role played by the suspect/witness biases $b_s,b_w$ in earlier
work: regularity of $\Gamma$ --- positive $p^\Gamma$ on every legal configuration
--- corresponds to the old $-0.5<b_s,b_w<0.5$ regularity. We will say that $\Gamma$
is \emph{regular} when $p^\Gamma(\mathbf{C}',\mathbf{T}')>0$ for every
$(\mathbf{C}',\mathbf{T}')$ with $\mathbf{C}'\subseteq\mathbf{S}$,
$\mathbf{T}'\subseteq\mathbf{W}$, $\mathbf{C}'\cap\mathbf{T}'=\emptyset$.

We can now record the forward direction.

\begin{theorem}\label{th:ac-exp}
    Suppose that $\mathbf{C}$ is an actual cause of $\varphi$ with witness set $\mathbf{T}$ in $(M,\mathbf{u})$. Take any $\{\mathbf{S}=\mathbf{s}\}\supseteq\mathbf{C}$ as the potential cause set, subject to the constraint that no element of $\mathbf{S}$ is a deterministic function (under $M$) of any subset of the remaining elements of $\mathbf{S}$,\footnote{Equivalently, $\mathbf{S}$ is closed under no nontrivial structural equation of $M$. The constraint excludes admitting both $S$ and a deterministic descendant $H=f(S,\ldots)$ as candidates in the same suspect set, but covers all candidate sets standardly considered in the AC literature, where structural descendants of one candidate are not admitted as separate candidates: their contribution is mediated by the equation, and AC's minimality clause already prefers the ancestor alone. The same closure condition discharges a joint-support concern: the next assumption requires $\Delta(\mathbf{s})$ to place positive density on every \emph{joint} alternative configuration $\mathbf{s}'\in\dom(\mathbf{S})\setminus\{\mathbf{s}\}$, not merely positive marginal density on each component. In a chain $S\to H$ with $H=f(S)$, joint assignments such as $(S{=}0, H{=}1)$ are impossible under the structural equations, so any $\Delta$ that placed mass on them would be incompatible with $M$. Closure rules out exactly this case: when no element of $\mathbf{S}$ is a deterministic function of the remaining elements, the joint domain $\dom(\mathbf{S})$ contains no rectangles forced to zero measure by $M$, and a $\Delta$ with positive PMF or Lebesgue density on every $\mathbf{s}'\neq\mathbf{s}$ is consistent with $M$ on the suspect set. In the stone-throwing benchmark of Section~\ref{sec:ac_benchmark}, the suspect set is $\{A_i,B_i\}_i$ and the deterministic mediators $\{W^a_i,W^b_i\}_i$ live in the witness set, not the suspect set, exactly to keep the suspect set closure-respecting. Witnesses are pinned at their factual values and are not subject to a joint-support requirement on alternatives, so they need no analogous condition.} and any $\mathbf{W}\supseteq\mathbf{T}$ as the potential witness set. Assume $\Gamma$ is regular and that $\Delta(\mathbf{s})$ assigns positive density to every $\mathbf{s}'\in\dom(\mathbf{S})\setminus\{\mathbf{s}\}$ (positive PMF in the discrete case, positive Lebesgue density in the continuous case). In the continuous case, assume additionally:
    \begin{itemize}
    \item[(i)] the structural equations of $M$ downstream of $\mathbf{C}$ are continuous in
    $\mathbf{c}'$ under fixed $\mathbf{u}$;
    \item[(ii)] $\neg\varphi$ is open in $\dom(Y)$ (or, more weakly, the boundary of
    $\neg\varphi$ has measure zero under the pushforward of $Y$).
    \end{itemize}
    Then
    \begin{align}
    \Phi(\mathbf{C},\mathbf{T}) > 0,
    \label{eq:pect}
    \end{align}
    and therefore
    \[
    \sum_{\mathbf{T}'\subseteq\mathbf{W}} \Phi(\mathbf{C},\mathbf{T}')\;>\;0.
    \]
\end{theorem}

\begin{proof}
Fix $\mathbf{u}$, $\mathbf{c}'$, and $\mathbf{T}$ as witnesses to the satisfaction of
the conditions in Definition~\ref{def:ac}. The context-sensitive-necessity clause
forces $\mathbf{T}\cap\mathbf{C}=\emptyset$, since otherwise the joint intervention
$[\mathbf{C}\!\leftarrow\!\mathbf{c}',\,\mathbf{T}\!\leftarrow\!\mathbf{t}^\star]$ would
be ill-defined. It remains to show that each of the three factors in
$\Phi(\mathbf{C},\mathbf{T})$ is strictly positive.
\begin{itemize}
\item \emph{Selection weight.} By regularity of $\Gamma$,
$p^\Gamma(\mathbf{C},\mathbf{T})>0$.
\item \emph{Sufficiency indicator.} With exogenous noise $\mathbf{u}$ held constant,
factivity gives $(M,\mathbf{u})\vDash\varphi$. Restoring $\mathbf{C}$ to
$\mathbf{s}^\star\!\mid_{\mathbf{C}}$ and $\mathbf{T}$ to $\mathbf{w}^\star\!\mid_{\mathbf{T}}$
leaves the model agreeing with $(M,\mathbf{u})$ on every variable, so the indicator
evaluates to $1$.
\item \emph{Necessity integral.} Context-sensitive necessity provides a specific
$\mathbf{c}'$ witnessing
$(M,\mathbf{u})\vDash[\mathbf{C}\!\leftarrow\!\mathbf{c}',\mathbf{T}\!\leftarrow\!\mathbf{t}^\star]\neg\varphi$.
The closure condition on $\mathbf{S}$ and the positive-density assumption on
$\Delta(\mathbf{s})$ imply that $\mathbf{c}'$ lies in the support of
$\Delta_{\mathbf{C}}(\mathbf{s}^\star)$. We treat the two cases separately.
\emph{Discrete case}: the singleton $\{\mathbf{c}'\}$ has positive PMF
$\Delta_{\mathbf{C}}(\{\mathbf{c}'\})>0$ by hypothesis, so the integral is at least
$\Delta_{\mathbf{C}}(\{\mathbf{c}'\})>0$. \emph{Continuous case}: hypothesis~(i) gives
continuity of the intervention map $\mathbf{c}'\mapsto Y$ at fixed $\mathbf{u}$;
hypothesis~(ii) places $Y(\mathbf{c}')$ in the interior of $\neg\varphi$; pulling back
an open neighbourhood within $\neg\varphi$ yields an open neighbourhood of
$\mathbf{c}'$ on which the indicator equals $1$, and this neighbourhood has positive
Lebesgue measure and hence positive $\Delta_{\mathbf{C}}$-mass by the
positive-density hypothesis. In either case the integral is strictly positive.
\end{itemize}
The product of three strictly positive factors is strictly positive, giving
$\Phi(\mathbf{C},\mathbf{T})>0$. Summing over $\mathbf{T}'\subseteq\mathbf{W}$ adds
non-negative terms to a positive one, preserving positivity. The proof uses factivity
and context-sensitive necessity but not minimality: the same conclusion holds whenever
$\mathbf{C}$ satisfies factivity and context-sensitive necessity in $(M,\mathbf{u})$,
a fact we will use in Proposition~\ref{prop:filter}.
\end{proof}

\medskip\noindent\textit{Remark (when (i)--(ii) bite).} In settings with continuous
outcomes and an open event such as a threshold $\{Y>y_0\}$ (for instance the SIR
overshoot benchmark of Section~\ref{sec:sir_benchmark}), both hypotheses are
satisfied trivially. Binary-output models, where $Y$ takes values in $\{0,1\}$, satisfy
(ii) automatically (every subset of a discrete codomain is open) but may not satisfy
(i) when the binary output comes from thresholding a continuous score. In such cases
the proof falls back to the discrete branch, which only requires positive PMF on the
witnessing $\mathbf{c}'$.

The opposite direction follows by reading the proof in reverse on the necessity factor.

\begin{theorem}\label{th:exp-ac}
Under the same hypotheses on $\{\mathbf{S}=\mathbf{s}\}$, $\Gamma$, and $\Delta$ as
in Theorem~\ref{th:ac-exp}, if
\[
\sum_{\mathbf{T}'\subseteq\mathbf{W}} \Phi(\mathbf{C},\mathbf{T}')\;>\;0,
\]
then both factivity ($(M,\mathbf{u})\vDash\varphi$) and the
context-sensitive-necessity condition in Definition~\ref{def:ac} hold for
$\mathbf{C}$ in $(M,\mathbf{u})$.
\end{theorem}

\begin{proof}
Since $\Phi$ is non-negative, the assumption forces $\Phi(\mathbf{C},\mathbf{T}')>0$
for at least one $\mathbf{T}'\subseteq\mathbf{W}$. By Definition~\ref{def:ackernel},
each factor of $\Phi(\mathbf{C},\mathbf{T}')$ is strictly positive. The sufficiency
indicator equals $1$, i.e.\
$(M,\mathbf{u})\vDash[\mathbf{C}\!\leftarrow\!\mathbf{s}^\star\!\mid_{\mathbf{C}},\,\mathbf{T}'\!\leftarrow\!\mathbf{w}^\star\!\mid_{\mathbf{T}'}]\varphi$;
since pinning variables to their factual values under fixed $\mathbf{u}$ leaves every
variable unchanged, this is equivalent to $(M,\mathbf{u})\vDash\varphi$, giving
factivity. The necessity integral being positive provides a
$\Delta_{\mathbf{C}}$-positive-density set of values $\mathbf{c}'$ for which
$(M,\mathbf{u})\vDash[\mathbf{C}\!\leftarrow\!\mathbf{c}',\,\mathbf{T}'\!\leftarrow\!\mathbf{w}^\star\!\mid_{\mathbf{T}'}]\neg\varphi$;
any such $\mathbf{c}'$, together with $\mathbf{T}'$, witnesses the existential in the
context-sensitive-necessity clause of Definition~\ref{def:ac}.
\end{proof}

\medskip\noindent\textit{Remark (deterministic intermediates and joint support).}
The hypothesis on $\{\mathbf{S}=\mathbf{s}\}$ is what rules out the following obstruction.
Suppose the suspect set were closed under the structural equations of $M$, e.g., both
Sally's throw $S$ and the bottle's hit $H = f(S,\ldots)$ admitted as candidates. Then the
joint assignment $(S{=}0, H{=}1)$ is structurally impossible, and any alternative-value
distribution consistent with $M$ assigns it density zero, so the joint full-support
condition required to recover the witnessing $\mathbf{c}'$ from $\Phi > 0$ fails.
Restricting suspects to variables that are not deterministic functions of the others
under $M$ removes this obstruction. The constraint is also conservative with respect to
the AC literature: in the canonical stone-throwing setup, the candidate set is
$\{S_{Sally}, S_{Bill}\}$, with $H$ as the outcome rather than a candidate, and AC's own
minimality clause already prefers ancestor candidates over their deterministic descendants
when both are admitted. Sections~\ref{sub:synthetic} and~\ref{sec:ac_benchmark} adopt the
same convention experimentally --- the continuous synthetic benchmark uses the six roots
as suspects and the discrete scaled-throwing benchmark uses $\{A_i,B_i\}_i$ as suspects
with the deterministic mediators $\{W^a_i,W^b_i\}_i$ as witnesses --- in both cases
deterministic intermediates are excluded from the suspect set. Cases where a user genuinely wants to attribute responsibility separately to a
deterministic intermediate (rather than to its structural ancestors) lie outside the
scope of these theorems and require an alternative reading of intervention semantics; we
treat that as a separate problem.

\medskip
With minimality, the situation is somewhat more complicated. Our choice of $\Gamma$
will have an impact on which configurations maximise $\Phi$. Moreover, $\Phi$ compares
configurations by $\Gamma$-weight rather than by the subset relation, while
Definition~\ref{def:ac} formulates minimality strictly in terms of subsets, so some
misalignment is to be expected. Some useful observations can be made nonetheless.

\begin{definition}[Necessity dilution]\label{def:necessity-dilution}
A pair $(M,\Delta)$ satisfies \emph{necessity dilution} on $\mathbf{S}$ if, writing
$N(\mathbf{C},\mathbf{T})$ for the necessity integral in
Equation~\eqref{eq:ackernel},
\[
N(\mathbf{C}^\sharp,\mathbf{T}^\sharp)\;\geq\;N(\mathbf{C}^\dagger,\mathbf{T}^\dagger)
\]
for every $\mathbf{C}^\sharp\subsetneq\mathbf{C}^\dagger\subseteq\mathbf{S}$ and
witnesses $\mathbf{T}^\sharp,\mathbf{T}^\dagger$ disjoint from the corresponding
suspect set.
\end{definition}

The first is that there is a simplified setting in which a positive value of $\Phi$,
under factivity, entails minimality with respect to the same parameters. Three
conditions suffice:
\begin{itemize}
\item $\Gamma_w$ is uniform over $2^{\mathbf{W}}$;
\item $\Gamma_s$ is \emph{strictly cardinality-decreasing}, i.e.\
$p^{\Gamma_s}(\mathbf{C}_1) > p^{\Gamma_s}(\mathbf{C}_2)$ whenever
$\mathbf{C}_1\subsetneq\mathbf{C}_2$ (the modern counterpart of the old null witness
bias and positive suspect preemption bias);
\item the pair $(M,\Delta)$ satisfies necessity dilution on $\mathbf{S}$
(Definition~\ref{def:necessity-dilution}).
\end{itemize}
The strategy used in the proof entails that we can safely marginalise over witnesses
as well.

\begin{theorem}\label{th:local_max_to_ac}
Consider the set of optimisers
$\mathcal{O}^\star = \arg\max_{\mathbf{C}\subseteq\{\mathbf{S}=\mathbf{s}\},\,\mathbf{T}\subseteq\mathbf{W}}
\Phi(\mathbf{C},\mathbf{T})$. Assume $(\mathbf{C}^\dagger,\mathbf{T}^\dagger)\in\mathcal{O}^\star$
with $\Phi(\mathbf{C}^\dagger,\mathbf{T}^\dagger)>0$, that factivity is satisfied for
$\mathbf{C}^\dagger$ (as in Definition~\ref{def:ac}), and that no element of
$\{\mathbf{S}=\mathbf{s}\}$ is a deterministic function under $M$ of the others. Let
$\Gamma_w$ be uniform on $2^{\mathbf{W}}$ and let $\Gamma_s$ be strictly
cardinality-decreasing, and assume that $(M,\Delta)$ satisfies necessity dilution on
$\mathbf{S}$. Then $\mathbf{C}^\dagger$ is an actual cause of $\varphi$ according to
Definition~\ref{def:ac}.
\end{theorem}

\begin{proof}
Factivity is assumed; by Theorem~\ref{th:exp-ac} the context-sensitive-necessity
condition holds for $\mathbf{C}^\dagger$ in $(M,\mathbf{u})$. So the only way
$\mathbf{C}^\dagger$ can fail to be an actual cause is by failing minimality. Suppose
for contradiction that minimality fails: there is
$\mathbf{C}^\sharp\subsetneq\mathbf{C}^\dagger$ satisfying factivity and
context-sensitive necessity for $\varphi$ with some witness $\mathbf{T}^\sharp$. By
Theorem~\ref{th:ac-exp} (its proof requires only factivity and necessity, not
minimality), $\Phi(\mathbf{C}^\sharp,\mathbf{T}^\sharp)>0$ as well, with the
sufficiency indicator and the necessity integral both equal to $1$ and a strictly
positive value, respectively. Both pairs therefore satisfy
$\Phi(\mathbf{C},\mathbf{T}) = p^\Gamma(\mathbf{C},\mathbf{T})\cdot N(\mathbf{C},\mathbf{T})$.

By the construction of $\Gamma$ from independent marginals
(Section~\ref{sec:causal_impact}),
\[
p^\Gamma(\mathbf{C},\mathbf{T}) \;\propto\;
p^{\Gamma_s}(\mathbf{C})\,p^{\Gamma_w}(\mathbf{T})\,
\mathbb{I}[\mathbf{T}\cap\mathbf{C}=\emptyset].
\]
The $\mathbb{I}[\cdot]$ factor is $1$ for both pairs (otherwise the contradiction
would have been immediate). Uniformity of $\Gamma_w$ on $2^{\mathbf{W}}$ makes
$p^{\Gamma_w}(\mathbf{T})=2^{-|\mathbf{W}|}$ identically, so the witness factor cancels.
Strict cardinality-decrease of $\Gamma_s$ together with
$\mathbf{C}^\sharp\subsetneq\mathbf{C}^\dagger$ gives
$p^{\Gamma_s}(\mathbf{C}^\sharp)>p^{\Gamma_s}(\mathbf{C}^\dagger)$, and necessity
dilution gives $N(\mathbf{C}^\sharp,\mathbf{T}^\sharp)\geq
N(\mathbf{C}^\dagger,\mathbf{T}^\dagger)$. Multiplying these (and noting
$N(\mathbf{C}^\sharp,\mathbf{T}^\sharp)>0$ to convert weak into strict in the
composite) gives
$\Phi(\mathbf{C}^\sharp,\mathbf{T}^\sharp)>\Phi(\mathbf{C}^\dagger,\mathbf{T}^\dagger)$,
contradicting $(\mathbf{C}^\dagger,\mathbf{T}^\dagger)\in\mathcal{O}^\star$.
\end{proof}

\begin{corollary}\label{cor:marginalize}
Keep the notation and hypotheses of Theorem~\ref{th:local_max_to_ac}. Consider the set
of optimisers
$\mathcal{O}^\Sigma = \arg\max_{\mathbf{C}\subseteq\{\mathbf{S}=\mathbf{s}\}}
\sum_{\mathbf{T}\subseteq\mathbf{W}} \Phi(\mathbf{C},\mathbf{T})$. Assume
$\mathbf{C}^\dagger\in\mathcal{O}^\Sigma$ with
$\sum_{\mathbf{T}\subseteq\mathbf{W}} \Phi(\mathbf{C}^\dagger,\mathbf{T})>0$ and that
factivity is satisfied for $\mathbf{C}^\dagger$. Then $\mathbf{C}^\dagger$ is an actual
cause of $\varphi$ according to Definition~\ref{def:ac}.
\end{corollary}

\begin{proof}
The summation hypothesis entails the existence of some
$\mathbf{T}^\dagger\subseteq\mathbf{W}$ with
$\Phi(\mathbf{C}^\dagger,\mathbf{T}^\dagger)>0$. Because the slice
$\{(\mathbf{C}^\dagger,\mathbf{T})\,:\,\mathbf{T}\subseteq\mathbf{W}\}$ is contained in
the joint search space $2^{\mathbf{S}}\times 2^{\mathbf{W}}$, we have
$(\mathbf{C}^\dagger,\mathbf{T}^\dagger)\in\mathcal{O}^\star$ as well. The conditions of
Theorem~\ref{th:local_max_to_ac} are satisfied, so $\mathbf{C}^\dagger$ is an actual
cause of $\varphi$.
\end{proof}

\medskip\noindent\textit{When necessity dilution holds.} The hypothesis is a property
of the pair $(M,\Delta)$, not of $\Gamma$. It says that adding suspects to a candidate
cause does not introduce new $\Delta$-mass of flipping alternatives, and three settings
make it automatic. \emph{(i) Single-pivot models}: one variable $X^\star\in\mathbf{S}$
is pivotal under $(M,\mathbf{u})$ and the others are structurally idle, so $N$ peaks at
$\{X^\star\}$ and falls off as more suspects are added. \emph{(ii) Factored $\Delta$
without Boolean-OR redundancy}: $\Delta_{\mathbf{C}}$ is a product of marginals
$\bigotimes_{X\in\mathbf{C}}\nu_X$ with $\nu_X(\dom(X))\leq 1$, and the flip event
factorises coordinate-wise; $N(\mathbf{C})$ is then a product of factors each at most
$1$, decreasing in $|\mathbf{C}|$. \emph{(iii) Conjunctive outcome events}: if
$\neg\varphi$ requires flipping every variable in a required set
$\mathbf{R}\subseteq\mathbf{S}$, the flip region's $\Delta_{\mathbf{C}}$-mass decreases
as $\mathbf{C}$ grows beyond $\mathbf{R}$. The hypothesis fails in disjunctive
overdetermination (e.g.\ $\varphi:\,A\vee B=1$ with both $A,B$ factual): no proper
subset of $\mathbf{S}$ flips $\varphi$, so $N$ is zero on subsets and positive on the
full set: the opposite of dilution. In such cases no proper subset of $\mathbf{S}$
is an actual cause anyway, so Theorem~\ref{th:local_max_to_ac} is not the appropriate
tool; Proposition~\ref{prop:filter} is.

\medskip\noindent We can now return to the misalignment we mentioned earlier. One can
imagine a reasoning mode in which one looks for causes of an outcome by inspecting
the integrated probability $P(\mathbf{C}\!\rightsquigarrow\!\varphi):=
\int\sum_{\mathbf{T}}\Phi(\mathbf{C},\mathbf{T})\,P_{\mathbf{U}}(d\mathbf{u})$, not knowing
what the values of $\mathbf{u}$ are but knowing that $\varphi$ occurred. It is worth
pointing out that
$\arg\max_{\mathbf{C}\subseteq\mathbf{S}} P(\mathbf{C}\!\rightsquigarrow\!\varphi)$
might fail to include some $\mathbf{C}$ that is an actual cause of $\varphi$ for some
setting $\mathbf{u}$.

\begin{observation}\label{obs:misalignment}
Define $\Xi^{\max}=\arg\max_{\mathbf{C}\subseteq\mathbf{S}} P(\mathbf{C}\!\rightsquigarrow\!\varphi)$.
It is not the case that if $\mathbf{C}^\dagger$ is an actual cause in $(M,\mathbf{u})$
for some $\mathbf{u}$, $\mathbf{C}^\dagger\in\Xi^{\max}$.
\end{observation}

\begin{proof}
Consider a model with two binary parents $A,B$ uniformly distributed and a binary
child $Y=A\vee B$. Take $\varphi:\,Y=1$ and the suspect set $\mathbf{S}=\{A,B\}$ with
$\mathbf{W}=\emptyset$. Compare $\mathbf{C}^\dagger=\{A,B\}$ at noise $\mathbf{u}_{11}$
($A=B=1$) and $\mathbf{C}^\sharp=\{A\}$ at noise $\mathbf{u}_{10}$ ($A=1,B=0$).
\begin{itemize}
\item In $\mathbf{u}_{11}$, $\mathbf{C}^\dagger$ is an actual cause: setting both
parents to $0$ yields $Y=0$, no proper subset achieves this, and the indicators in
$\Phi$ are $1$. With $\Gamma_w$ uniform on $2^{\emptyset}$ (a single configuration of
weight $1$) and a strictly cardinality-decreasing $\Gamma_s$,
$\Phi(\mathbf{C}^\dagger,\emptyset)\propto p^{\Gamma_s}(\{A,B\})$.
\item In $\mathbf{u}_{10}$, $\mathbf{C}^\sharp$ is an actual cause:
intervening to $A=0$ flips $Y$ to $0$, the singleton is minimal, and again the
indicators are $1$. So $\Phi(\mathbf{C}^\sharp,\emptyset)\propto p^{\Gamma_s}(\{A\})$.
\end{itemize}
Integrating uniformly over the four noise realisations $\{\mathbf{u}_{ij}\}$,
$P(\mathbf{C}^\dagger\!\rightsquigarrow\!\varphi)\propto p^{\Gamma_s}(\{A,B\})$ (only
$\mathbf{u}_{11}$ contributes) while
$P(\mathbf{C}^\sharp\!\rightsquigarrow\!\varphi)\propto 2\,p^{\Gamma_s}(\{A\})$
($\mathbf{u}_{10}$ and $\mathbf{u}_{11}$ both contribute). Whether $\mathbf{C}^\dagger$
or $\mathbf{C}^\sharp$ ends up in $\Xi^{\max}$ depends on the precise shape of
$\Gamma_s$: under sharply decreasing $\Gamma_s$, $\mathbf{C}^\sharp$ wins; under
shallow decrease, $\mathbf{C}^\dagger$ may. Both cannot belong to $\Xi^{\max}$
simultaneously in general, so some actual causes are dropped.
\end{proof}

The reason is structural: actual cause sets may be subset-minimal without being
cardinality-minimal, while integrated-probability maximisation is governed by
$\Gamma_s$, which (under the cardinality-decreasing parameterisation that aligns with
AC's minimality) penalises larger sets. The misalignment is therefore not a defect of
$\Phi$ but a mismatch between AC's notion of minimality (subset relation) and
$\Gamma$-weighted scoring.

The fix is to drop the maximisation step and use $\Phi$ directly to enumerate
candidate causes, then filter for subset-minimality. The filter recovers exactly the
actual causes in $\mathbf{S}$.

\begin{proposition}\label{prop:filter}
Under the hypotheses of Theorem~\ref{th:ac-exp}, define
\[
\Xi \;:=\; \Big\{\mathbf{C}^\dagger\subseteq\mathbf{S}\;:\;
\sum_{\mathbf{T}\subseteq\mathbf{W}} \Phi(\mathbf{C}^\dagger,\mathbf{T})>0,\;
\nexists\,\mathbf{C}'\subsetneq\mathbf{C}^\dagger\text{ with }
\sum_{\mathbf{T}\subseteq\mathbf{W}} \Phi(\mathbf{C}',\mathbf{T})>0
\Big\}.
\]
Then $\Xi$ contains exactly the actual causes of $\varphi$ in $(M,\mathbf{u})$ that are
subsets of $\mathbf{S}$.
\end{proposition}

\begin{proof}
\emph{$\Xi$-membership $\Rightarrow$ actual cause.} Let $\mathbf{C}^\dagger\in\Xi$.
Theorem~\ref{th:exp-ac} applied to $\sum_{\mathbf{T}}\Phi(\mathbf{C}^\dagger,\mathbf{T})>0$
gives both factivity and context-sensitive necessity. For minimality, suppose
$\mathbf{C}_{\text{sub}}\subsetneq\mathbf{C}^\dagger$ satisfied factivity and
context-sensitive necessity. By the proof of Theorem~\ref{th:ac-exp} (which uses
only factivity and necessity, not minimality), we would have
$\sum_{\mathbf{T}}\Phi(\mathbf{C}_{\text{sub}},\mathbf{T})>0$, contradicting the
no-positive-proper-subset clause of $\Xi$.

\emph{Actual cause $\Rightarrow$ $\Xi$-membership.} Let $\mathbf{C}^\dagger\subseteq\mathbf{S}$
be an actual cause of $\varphi$ in $(M,\mathbf{u})$. Theorem~\ref{th:ac-exp} gives
$\sum_{\mathbf{T}}\Phi(\mathbf{C}^\dagger,\mathbf{T})>0$. For the no-positive-proper-subset
clause: suppose $\mathbf{C}'\subsetneq\mathbf{C}^\dagger$ had
$\sum_{\mathbf{T}}\Phi(\mathbf{C}',\mathbf{T})>0$. Then by Theorem~\ref{th:exp-ac},
both factivity and context-sensitive necessity would hold for $\mathbf{C}'$. So
$\mathbf{C}'$ would witness a violation of $\mathbf{C}^\dagger$'s minimality,
contradicting the assumption that $\mathbf{C}^\dagger$ is an actual cause.
\end{proof}

\noindent Proposition~\ref{prop:filter} converts the cardinality--subset misalignment
into a procedure: enumerate $\{\mathbf{C}\subseteq\mathbf{S}:\sum_{\mathbf{T}}\Phi(\mathbf{C},\mathbf{T})>0\}$
(which by Theorem~\ref{th:exp-ac} are exactly the cause sets satisfying
context-sensitive necessity), filter for factivity and subset-minimality, and read off
the actual causes. Some idiosyncratic features of the AC notion (in particular, the
subset-minimality clause, which interacts poorly with $\Gamma$-weighted scoring) are
part of what motivated the more scalable \ourapproach{} formulation of
Section~\ref{sec:causal_impact} in the first place; the filter procedure here is the
right tool when AC verdicts specifically are wanted, but expectation-based
\ourapproach{} attributions are what scale to ML-grade models.

\medskip\noindent\textit{Summary of guarantees.} The section delivers, under explicit
hypotheses on $\mathbf{S}$, $\Gamma$, and $\Delta$:
(i) a forward direction (Theorem~\ref{th:ac-exp})
``actual cause $\Rightarrow\Phi>0$'';
(ii) a backward direction (Theorem~\ref{th:exp-ac})
``$\sum_{\mathbf{T}'}\Phi>0\Rightarrow$ context-sensitive necessity'';
(iii) minimality preservation under cardinality-decreasing $\Gamma_s$
(Theorem~\ref{th:local_max_to_ac}, Corollary~\ref{cor:marginalize});
(iv) a counterexample (Observation~\ref{obs:misalignment}) showing that integration
over noise misaligns with subset-minimality; and
(v) a corrective filter procedure (Proposition~\ref{prop:filter}) that recovers AC
verdicts exactly. The empirical counterpart of the correspondence, a benchmark of
the necessity-only PCI estimator against exact subset enumeration on a scaled
stone-throwing problem, is reported in Section~\ref{sec:ac_benchmark}.

\paragraph{Where we go next.} Before turning to the empirical benchmark, the next
section locates \ourapproach against the construction in the literature closest to
its expectation form: Pearl's \emph{probability of actual causation}. The two share
the same surface idea, ``how often, under the noise, does the actual-cause
verdict fire?'', and they coincide on standard examples, but they disagree on
graded responsibility patterns that PCI registers and Pearl's binary AC verdict
discards. Section~\ref{sub:pearl_actual_cause_prob} walks through Pearl's
canonical desert-traveller illustration and a weak-poison variant to make the
distinction concrete.
